% =====================================================================
%  Chapter 9: Spatial Dependence and Autocorrelation
%  Source: SpatialStats.tex (motivations prose) + P4 transcription (full)
% =====================================================================

\chapter{Spatial Dependence and Autocorrelation}
\label{chap_Spatial}

\section{Motivations for Going Spatial}
\label{sec_MotivationsSpatial}

Spatial autocorrelation\index{spatial autocorrelation} can be loosely defined
as ``the coincidence of value similarity with locational similarity''
\citep{Anselin2002}. The problem of spatial dependence is best viewed as a
normal extension of Tobler's First Law of Geography:

\begin{quote}
``everything is related to everything else, but near things are more related
than distant things'' \citep[p.236]{Tobler1970}.
\end{quote}

Three main reasons motivate the use of spatial methods:

\begin{enumerate}[label=(\roman*)]
  \item \textbf{The independence assumption is not valid:} Attributes of
    observation $i$ may influence attributes of observation $j$.
  \item \textbf{Spatial heterogeneity:} The magnitude and direction of the
    treatment effect may vary across space.
  \item \textbf{Omitted variable bias:} Some common unobserved or latent
    variables are shared among spatial neighbours.
\end{enumerate}

Recall the non-spatial DGP:
\begin{equation}
  y_i = X_i\beta + e_i, \qquad e_i \sim \text{iid}(0,\sigma^2)
  \label{eq_NonSpatDGP}
\end{equation}

With assumptions: (1) observed values at location $i$ are independent of
values at location $j$; (2) $\E[e_ie_j] = \E[e_i]\E[e_j] = 0$.

\begin{tcolorbox}[definition]
\textbf{MAUP Warning.} The choice of spatial weight matrix is not merely
technical---it is an inferential decision with the same logical status as the
choice of instruments in IV estimation. Results should always be checked for
sensitivity to weight matrix specification.
\end{tcolorbox}


\section{The Spatial DGP}
\label{sec_SpatialDGP}

With 2 neighbours $i$ and $j$, observations depend on each other:
\begin{alignat*}{2}
  y_i &= \alpha_j y_j + X_i\beta + e_i, \qquad &e_i &\sim \text{iid}(0,\sigma^2) \\
  y_j &= \alpha_i y_i + X_j\beta + e_j, \qquad &e_j &\sim \text{iid}(0,\sigma^2)
\end{alignat*}

\noindent\textbf{Problem.} With more spatial dependencies, there is little
practical usefulness since the system has \emph{more parameters than
observations}: potentially $n^2 - n$ relations to be estimated.

\noindent\textbf{Assumptions of the spatial DGP:}
\begin{itemize}
  \item Observations at location $i$ depend on those at location $j$ and
    vice versa.
  \item The DGP is \emph{simultaneous} across space.
\end{itemize}

The overparametrisation problem is solved by imposing structure in a
\textbf{spatial autoregressive process}\index{spatial autoregressive process}.
Introducing a scalar spatial correlation coefficient $\rho$:

\begin{tcolorbox}[keyresult]
\begin{equation}
  y_i = \rho\sum_{j=1}^{n} W_{ij}\,y_j + e_i
\end{equation}
where $\sum_j W_{ij}y_j$ is the \textbf{spatial lag}\index{spatial lag} of
$y_i$---the weighted average of neighbouring observations.
\end{tcolorbox}


\section{The Spatial Weight Matrix}
\label{sec_WeightMatrix}

In matrix notation ($W$ of size $n\times n$; $y$ of size $n\times 1$):
\begin{equation}
  y = \rho Wy + e, \qquad e \sim N(0, \sigma^2 I_n)
\end{equation}

To construct $W$, begin with a \textbf{contiguity matrix}
$C$\index{contiguity matrix}, where rows represent observation $i$ and
columns represent neighbour $j$ (1 = neighbour, 0 = not).

\subsection{Example: Four Contiguous Regions}

Consider four regions $R_1$--$R_4$ arranged in a chain. First-order
neighbours ($N_1$) share a boundary; $R_1$ is its own second-order
neighbour (the neighbour of its first-order neighbour $R_2$).

\begin{figure}[htp]
\centering
\begin{tikzpicture}[>=Stealth, node distance=0cm]
  \node[draw, minimum width=1.2cm, minimum height=0.8cm] (R1) {$R_1$};
  \node[draw, minimum width=1.2cm, minimum height=0.8cm, right=of R1] (R2) {$R_2$};
  \node[draw, minimum width=1.2cm, minimum height=0.8cm, right=of R2] (R3) {$R_3$};
  \node[draw, minimum width=1.2cm, minimum height=0.8cm, right=of R3] (R4) {$R_4$};
  \draw[<->, thick] ([yshift=6pt]R1.north) to[bend left=40]
    node[above] {$N_1$} ([yshift=6pt]R2.north);
  \draw[<->, thick] ([yshift=-6pt]R2.south) to[bend right=40]
    node[below] {$N_2$} ([yshift=-6pt]R3.south);
  \draw[<->, thick] ([yshift=-6pt]R3.south) to[bend right=40]
    node[below] {$N_2$} ([yshift=-6pt]R4.south);
\end{tikzpicture}
\caption{Four contiguous regions with first-order ($N_1$) and
higher-order ($N_2$) neighbours.}
\label{fig_ContiguousRegions}
\end{figure}

The contiguity matrix:
\begin{equation}
  C = \begin{bmatrix}
    0 & 1 & 0 & 0 \\
    1 & 0 & 1 & 0 \\
    0 & 1 & 0 & 1 \\
    0 & 0 & 1 & 0
  \end{bmatrix}
\end{equation}

Create $W$ by \textbf{row-standardisation} (row sums normalised to one):
\begin{equation}
  W = \begin{bmatrix}
    0   & 1   & 0   & 0   \\
    1/2 & 0   & 1/2 & 0   \\
    0   & 1/2 & 0   & 1/2 \\
    0   & 0   & 1   & 0
  \end{bmatrix}
\end{equation}

\noindent\textbf{Verifying the spatial lag $Wy$:}
\begin{equation}
  Wy = \begin{bmatrix}y_2\\(y_1+y_3)/2\\(y_2+y_4)/2\\y_3\end{bmatrix}
\end{equation}

Each element of $Wy$ is the (row-weighted) average of the neighbours of the
corresponding observation.

\subsection{Weight Matrix Types}

\begin{itemize}
  \item \textbf{Rook contiguity:} Shares a boundary (no diagonal touches)
  \item \textbf{Queen contiguity:} Shares a boundary or corner
  \item \textbf{$k$-nearest neighbours:} Fixed number of nearest units
  \item \textbf{Distance decay:} $w_{ij} = d_{ij}^{-\delta}$ for some
    $\delta > 0$, often with a distance cutoff
  \item \textbf{Row-standardised:} Rows sum to one; enables spatial lag
    interpretation as a local average
\end{itemize}


\section{The Spatial Autoregressive (SAR) DGP}
\label{sec_SARDGP}

$\rho$ is a scalar parameter of spatial dependence:
\begin{equation}
  \rho \;\begin{cases}
    = 0 & \text{no spatial autocorrelation} \\
    > 0 & \text{positive spatial autocorrelation (clustering)} \\
    < 0 & \text{negative spatial autocorrelation (segregation)}
  \end{cases}
\end{equation}

The full \textbf{SA DGP} with constant term and covariates:
\begin{align}
  y &= \rho Wy + \boldsymbol{\iota}\alpha + e \\
  (I_n - \rho W)y &= \boldsymbol{\iota}\alpha + e \\
  y &= (I_n - \rho W)^{-1}\boldsymbol{\iota}\alpha + (I_n - \rho W)^{-1}e
\end{align}

Assuming $|\rho| < 1$ (stationarity condition; see Section~\ref{sec_Eigenvalues}
for the general condition $|\rho| < 1/\lambda_{\max}(W)$), the inverse can be
expressed as a \textbf{geometric matrix series}:
\begin{equation}
  (I_n - \rho W)^{-1} = I_n + \rho W + \rho^2 W^2 + \rho^3 W^3 + \cdots
\end{equation}

Therefore:
\begin{equation}
  y = \boldsymbol{\iota}\alpha + \rho W\boldsymbol{\iota}\alpha
      + \rho^2 W^2\boldsymbol{\iota}\alpha + \cdots
      + e + \rho We + \rho^2 W^2 e + \cdots
\end{equation}

The first part converges to $(1-\rho)^{-1}\boldsymbol{\iota}\alpha$ (since
$\alpha$ is scalar and $W\boldsymbol{\iota} = \boldsymbol{\iota}$):

\begin{tcolorbox}[keyresult]
\begin{equation}
  y = \frac{1}{1-\rho}\,\boldsymbol{\iota}\alpha
      + \underbrace{e + \rho We + \rho^2 W^2 e + \cdots}_{\text{violation of Gauss-Markov}}
\end{equation}
\end{tcolorbox}

This is a \textbf{strong violation} of the Gauss-Markov assumption of
independent error terms. Therefore \textbf{OLS is biased and inconsistent}
under a SAR DGP.

\subsection{SAR Model with Covariates}

\begin{equation}
  y = \rho Wy + X\beta + e
  \;\Longrightarrow\;
  y = (I_n-\rho W)^{-1}X\beta + (I-\rho W)^{-1}e, \quad e\sim N(0,\sigma^2I)
\end{equation}

Estimate three parameters: $\beta$, $\sigma^2$, and $\rho$.
If $\rho=0$, the SAR model reduces to standard OLS.

The SAR estimator can be written as:

\begin{tcolorbox}[keyresult]
\begin{equation}
  \hat\beta_\text{SAR} = \hat\beta_\text{OLS} - \hat\rho(X'X)^{-1}X'Wy
  \label{eq_SARbeta}
\end{equation}
\end{tcolorbox}


\section{Global Spatial Autocorrelation: Moran's $I$}
\label{sec_MoranI}

Moran's $I$ coefficient calculates the ratio of the product of a variable
and its spatial lag, adjusted for the sum of spatial weights, relative to
the overall variance:

\begin{tcolorbox}[keyresult]
\begin{equation}
  I = \frac{\displaystyle
       \frac{\displaystyle\sum_i\sum_j w_{ij}(y_i-\bar{y})(y_j-\bar{y})}
            {\displaystyle\sum_i\sum_j w_{ij}}}
      {\displaystyle\frac{\sum_i(y_i-\bar{y})^2}{n}}
  \;\in\; [-1,\,+1]
  \label{eq_MoranI}
\end{equation}
\end{tcolorbox}

$I = +1$: perfect spatial clustering; $I = -1$: perfect dispersion;
$I \approx 0$: spatial randomness. Moran's $I$ can equivalently be
interpreted as the \textbf{regression coefficient of a variable on its
spatial lag}---the slope of the Moran scatterplot.

\subsection{Local Moran's $I$}

The global statistic can be decomposed into contributions from individual
observations:
\begin{equation}
  I_i = \frac{(y_i-\bar{y})\displaystyle\sum_j w_{ij}(y_j-\bar{y})}
             {\displaystyle\frac{\sum_i(y_i-\bar{y})^2}{n}}
  \label{eq_LocalMoran}
\end{equation}

Local Moran's $I$ \citep{Anselin1995} is the basis of LISA (Local Indicators
of Spatial Association) analysis. A map of local $I_i$ $z$-scores identifies:
\begin{itemize}
  \item \textbf{HH} (high-high): high-value location surrounded by
    high-value neighbours---core of a spatial cluster
  \item \textbf{LL} (low-low): low-value surrounded by low-value neighbours
  \item \textbf{HL/LH} (high-low or low-high): spatial outliers---locations
    significantly different from their neighbourhood
\end{itemize}


\section{Estimating Models with a SAR DGP}
\label{sec_EstimatingSAR}

Recall: OLS is inconsistent and biased under a SAR DGP. Use alternative
estimators: \textbf{MLE} or \textbf{IV/2SLS}. 2SLS has similar properties
to MLE (consistency and asymptotic normality) but is considerably easier
to implement numerically.

\subsection{SAR-IV Estimator}

Address endogeneity due to spatial variable linkage by using lagged variables
as instruments. Let $X = [Wy\;\; X_1]$ and $Z = [L(y)\;\; X_1]$, where
$L(y)$ denotes higher-order spatial lags. Then:
\begin{equation}
  \hat\beta_\text{SAR-IV} = [X'P_ZX]^{-1}X'P_Zy,
  \qquad P_Z \equiv Z[Z'Z]^{-1}Z'
\end{equation}


\section{The Spatial Model Taxonomy}
\label{sec_SpatialTaxonomy}

All models below are expressed relative to OLS ($y = X\beta + e$):

\renewcommand{\arraystretch}{1.5}
\begin{table}[htp]
\centering
\caption{Spatial regression model taxonomy}
\label{tbl_SpatialModels}
\begin{tabular}{@{}lp{5.5cm}p{4.5cm}@{}}
\toprule
\textbf{Model} & \textbf{Equation} & \textbf{Notes} \\
\midrule
OLS  & $y = X\beta + e$ & Baseline; assumes spatial independence \\
SAR  & $y = \rho Wy + X\beta + e$ & Spatial lag in $y$; OLS biased \\
SEM  & $y = X\beta + (I-\lambda W)^{-1}e$ & Spatial error model; OLS unbiased but inefficient \\
SDM  & $y = \rho Wy + X\beta + WX\theta + e$ & Spatial Durbin model \\
SARMA & $y = \rho Wy + X\beta + (I-\theta W)^{-1}e$ & SAR + MA error \\
SAC  & $y = \rho W_1y + X\beta + (I-\lambda W_2)^{-1}e$ & Spatial dependence in $y$ and $e$; separate weight matrices \\
SDEM & $y = X\beta + WX\theta + (I-\lambda W)^{-1}e$ & Spatial Durbin error model \\
\bottomrule
\end{tabular}
\end{table}
\renewcommand{\arraystretch}{1}


\section{Geographically Weighted Regression}
\label{sec_GWR}

Standard OLS assumes $\beta$ is globally constant. GWR relaxes this by
allowing locally-varying parameter estimates:
\begin{equation}
  y_i = X\beta_i + e \qquad\text{with locally-varying }\beta_i
\end{equation}

\begin{tcolorbox}[keyresult]
\begin{equation}
  \hat\beta_\text{GWR}(i) = (X'W_iX)^{-1}X'W_iy
  \label{eq_GWR}
\end{equation}
\end{tcolorbox}

Weights $w_{ij}$ are calculated using a \textbf{Gaussian decay
kernel}\index{Gaussian kernel!GWR}:
\begin{equation}
  w_{ij} = e^{(-d_{ij}/h)^2}
\end{equation}
where $d_{ij}$ is the Euclidean distance between observations $i$ and $j$,
and $h$ is the user-selected \textbf{bandwidth}\index{bandwidth!GWR}
(determined by cross-validation).

\noindent\textbf{Interpreting spatial variation in coefficients.}
GWR maps of $\hat\beta_i$ show where the relationship between $X$ and $y$
is stronger or weaker. However, spatial variation in coefficients can reflect
either genuine spatial heterogeneity in the relationship or noise arising from
sparse data. Significance tests for GWR coefficients should be interpreted
cautiously, as they rely on asymptotic approximations that may not be valid
for small local sample sizes.


\section{A Brief Excursion on Map Projections}
\label{sec_Projections_chap9}

\subsection{Coordinate Systems}

\noindent\textbf{Lat:} North--south position $\equiv$ $y$-axis\hfill
\noindent\textbf{Lon:} East--west position $\equiv$ $x$-axis

Projection is necessary to represent the 3D globe on a 2D map; projections
always distort some property of the original globe.

\begin{figure}[htp]
\centering
\begin{tikzpicture}[>=Stealth, font=\small]
  % Cylindrical
  \begin{scope}[xshift=0cm]
    \draw[thick] (-0.7,-1.1) -- (-0.7,1.1);
    \draw[thick] ( 0.7,-1.1) -- ( 0.7,1.1);
    \draw[thick] (-0.7,1.1) arc[start angle=180,end angle=0,x radius=0.7,y radius=0.18];
    \draw[thick,dashed] (-0.7,1.1) arc[start angle=180,end angle=360,x radius=0.7,y radius=0.18];
    \draw[thick] (-0.7,-1.1) arc[start angle=180,end angle=360,x radius=0.7,y radius=0.18];
    \draw[thick,dashed] (-0.7,-1.1) arc[start angle=180,end angle=0,x radius=0.7,y radius=0.18];
    \draw[thick] (0,0) circle (0.55);
    \node[below] at (0,-1.3) {(i) Cylindrical};
  \end{scope}
  % Conic
  \begin{scope}[xshift=3.5cm]
    \draw[thick] (0,1.5) -- (-0.85,-0.3);
    \draw[thick] (0,1.5) -- ( 0.85,-0.3);
    \draw[thick] (-0.85,-0.3) arc[start angle=180,end angle=360,x radius=0.85,y radius=0.2];
    \draw[thick,dashed] (-0.85,-0.3) arc[start angle=180,end angle=0,x radius=0.85,y radius=0.2];
    \draw[thick] (0,-0.1) circle (0.55);
    \node[below] at (0,-1.0) {(ii) Conic};
  \end{scope}
  % Planar
  \begin{scope}[xshift=7.0cm]
    \fill[gray!15] (-0.9,0.2) -- (0,1.2) -- (0.9,0.2) -- (0,-0.8) -- cycle;
    \draw[thick]   (-0.9,0.2) -- (0,1.2) -- (0.9,0.2) -- (0,-0.8) -- cycle;
    \draw[thick] (0,0) circle (0.5);
    \node[below] at (0,-1.0) {(iii) Planar};
  \end{scope}
\end{tikzpicture}
\caption{Three main types of map projection.}
\label{fig_Projections}
\end{figure}

\subsection{Projections and What They Preserve}

\renewcommand{\arraystretch}{1.4}
\begin{table}[htp]
\centering
\caption{Projection names by preserved property}
\label{tbl_Projections}
\begin{tabular}{@{}llll@{}}
\toprule
\textbf{Preserves} & \textbf{Cylindrical} & \textbf{Conic} & \textbf{Planar} \\
\midrule
Area     & \texttt{cylegnal\_area} & \texttt{albers}  & \texttt{azequal\_area} \\
Angle    & \texttt{mercator} (UTM, omerc) & \texttt{lambert} & \texttt{stereographic} \\
Distance & ---                    & ---              & \texttt{azequidistant}  \\
\bottomrule
\end{tabular}
\end{table}
\renewcommand{\arraystretch}{1}

\subsection{The Tissot Indicatrix}

The \textbf{Tissot indicatrix}\index{Tissot indicatrix} characterises
distortions due to map projections. Infinitesimally small circles of equal
radius placed at regular intervals on the sphere become ellipses after
projection; their size, shape, and orientation reveal how the projection
distorts area, angle, and distance.

\renewcommand\bibname{Further reading}
\begin{subappendices}
\bibliographystyle{econometrica}
\bibliography{Literature/OverLord}
\IfFileExists{Literature/SpatialStats.bbl}{\chapter{Applied Spatial Data Analysis}\label{chap_SpatialStats}
Quantitative analysis of spatial data has advanced greatly in recent years, largely as the result of the rapidly decreasing cost of computational power, the increased availability of high-resolution data and methodological innovation in spatial econometrics. Indeed, recent analytical and technological advances are increasingly blurring the traditional boundaries in spatial data analysis where various applied methods in Geographic Information Systems (GIS)\index{GIS|see{geographic information systems}}\index{geographic information systems}, geology-oriented remote sensing\index{remote sensing} and spatial statistics\index{spatial statistics} were essentially treated as separate fields. The legacy of this separate treatment of spatial data is reflected by the fact that even the leading textbooks in this area do not focus on more than two of these three sub-fields \citep[e.g.][]{Davis2002a,O'Sullivan2003,Campbell2006}.

Nowhere in the social sciences, perhaps, is this silent revolution felt more than in urban and regional planning. While planners usually have emphasised the application to GIS for land-use and environmental planning purposes, sometimes in addition with a little attention being paid to either remote sensing and spatial statistics, the discipline has never embraced all three sub-fields of quantitative data analysis in a systematic way. 

However, with the rapid deployment of ``Big Data''\index{Big Data} and the emergence of ``policy informatics''\index{policy informatics} as a new paradigm among computational social scientists, planners should increasingly see themselves at the forefront of the application of such integrated spatial data methods.\footnote{See a recent study by McKinsey \& Co. on this \href{http://www.mckinsey.com/Insights/MGI/Research/Technology\_and\_Innovation/Big\_data\_The\_next\_frontier\_for\_innovation}{topic}. See \citet{Barrett2011} for a definition of the term and on recent developments in the area of policy informatics.} \citet{Anselin2012} provides a good overview that assesses the evolution of the way in which spatial data analytical methods have been incorporated into software tools over the past two decades.


\section{Motivations for ``going spatial''}

Spatial autocorrelation\index{spatial autocorrelation|see{autocorrelation}}\index{autocorrelation!spatial} can be loosely defined as ``the coincidence of value similarity with locational similarity'' \citep{Anselin2002}. The problem of spatial dependence then is best viewed as a normal extension of the first law in geography, whereby

\begin{quote}
``everything is related to everything else, but near things are more related than distant things'' \citep[p.236]{Tobler1970}.
\end{quote}

Broadly speaking, there are three main reasons for deploying spatial methods:

\begin{enumerate}[(i)]
\item \emph{The independence assumption is not valid}: The attributes of observation $i$ may influence the attributes of observation
$j$.
\item \emph{Spatial heterogeneity}: The magnitude and direction the the treatment effect may vary across space.
\item \emph{Omitted variable bias}: Some common unobserved or latent variables are shared among spatial neighbours.
\end{enumerate}


Recall that the matrix notation of the non-spatial DGP was defined as

\begin{eqnarray}
\boldsymbol{y} = \boldsymbol{X\beta} + \boldsymbol{e} \qquad \text{with} \qquad \boldsymbol{e}\sim\mathcal{N}(0,\sigma^{2}\boldsymbol{I}).\label{eqn_SS_DGP}
\end{eqnarray}

Similarly, also recall that the classical normal regression model (CNR) of equation~\ref{eqn_SS_DGP} is based on the following assumptions:

\begin{itemize}
\item Observed values at location $i$ are independent of values at location $j$.
\item The regression residuals associated observation $i$ are independent of those associated with observation $j$, i.e. $\mathbb{E}[e_{i}e_{j}]=\mathbb{E}[e_{i}]\mathbb{E}[e_{j}]=0$.
\end{itemize}

\begin{itemize}
\item Include the following literature: \citet{Bivand2008}, \citet{Bivand2008a}, \citet{Corrado2012}, \citet{Kelejian1998}, \citet{McMillen2010a}, \citet{Pinske2010}.
\end{itemize}




\section{Spatial point pattern analysis}
\section{Interpolation and geostatistics}
\section{Areal data and spatial autocorrelation}

\begin{subappendices}
\section{Code for Spatial Stats in R}
\lstinputlisting{"Programs/UP504_SpatStat1.R"}

\renewcommand\bibname{Further reading}

\bibliographystyle{econometrica}
\bibliography{Literature/OverLord}
\IfFileExists{Literature/SpatialStats.bbl}{\chapter{Applied Spatial Data Analysis}\label{chap_SpatialStats}
Quantitative analysis of spatial data has advanced greatly in recent years, largely as the result of the rapidly decreasing cost of computational power, the increased availability of high-resolution data and methodological innovation in spatial econometrics. Indeed, recent analytical and technological advances are increasingly blurring the traditional boundaries in spatial data analysis where various applied methods in Geographic Information Systems (GIS)\index{GIS|see{geographic information systems}}\index{geographic information systems}, geology-oriented remote sensing\index{remote sensing} and spatial statistics\index{spatial statistics} were essentially treated as separate fields. The legacy of this separate treatment of spatial data is reflected by the fact that even the leading textbooks in this area do not focus on more than two of these three sub-fields \citep[e.g.][]{Davis2002a,O'Sullivan2003,Campbell2006}.

Nowhere in the social sciences, perhaps, is this silent revolution felt more than in urban and regional planning. While planners usually have emphasised the application to GIS for land-use and environmental planning purposes, sometimes in addition with a little attention being paid to either remote sensing and spatial statistics, the discipline has never embraced all three sub-fields of quantitative data analysis in a systematic way. 

However, with the rapid deployment of ``Big Data''\index{Big Data} and the emergence of ``policy informatics''\index{policy informatics} as a new paradigm among computational social scientists, planners should increasingly see themselves at the forefront of the application of such integrated spatial data methods.\footnote{See a recent study by McKinsey \& Co. on this \href{http://www.mckinsey.com/Insights/MGI/Research/Technology\_and\_Innovation/Big\_data\_The\_next\_frontier\_for\_innovation}{topic}. See \citet{Barrett2011} for a definition of the term and on recent developments in the area of policy informatics.} \citet{Anselin2012} provides a good overview that assesses the evolution of the way in which spatial data analytical methods have been incorporated into software tools over the past two decades.


\section{Motivations for ``going spatial''}

Spatial autocorrelation\index{spatial autocorrelation|see{autocorrelation}}\index{autocorrelation!spatial} can be loosely defined as ``the coincidence of value similarity with locational similarity'' \citep{Anselin2002}. The problem of spatial dependence then is best viewed as a normal extension of the first law in geography, whereby

\begin{quote}
``everything is related to everything else, but near things are more related than distant things'' \citep[p.236]{Tobler1970}.
\end{quote}

Broadly speaking, there are three main reasons for deploying spatial methods:

\begin{enumerate}[(i)]
\item \emph{The independence assumption is not valid}: The attributes of observation $i$ may influence the attributes of observation
$j$.
\item \emph{Spatial heterogeneity}: The magnitude and direction the the treatment effect may vary across space.
\item \emph{Omitted variable bias}: Some common unobserved or latent variables are shared among spatial neighbours.
\end{enumerate}


Recall that the matrix notation of the non-spatial DGP was defined as

\begin{eqnarray}
\boldsymbol{y} = \boldsymbol{X\beta} + \boldsymbol{e} \qquad \text{with} \qquad \boldsymbol{e}\sim\mathcal{N}(0,\sigma^{2}\boldsymbol{I}).\label{eqn_SS_DGP}
\end{eqnarray}

Similarly, also recall that the classical normal regression model (CNR) of equation~\ref{eqn_SS_DGP} is based on the following assumptions:

\begin{itemize}
\item Observed values at location $i$ are independent of values at location $j$.
\item The regression residuals associated observation $i$ are independent of those associated with observation $j$, i.e. $\mathbb{E}[e_{i}e_{j}]=\mathbb{E}[e_{i}]\mathbb{E}[e_{j}]=0$.
\end{itemize}

\begin{itemize}
\item Include the following literature: \citet{Bivand2008}, \citet{Bivand2008a}, \citet{Corrado2012}, \citet{Kelejian1998}, \citet{McMillen2010a}, \citet{Pinske2010}.
\end{itemize}




\section{Spatial point pattern analysis}
\section{Interpolation and geostatistics}
\section{Areal data and spatial autocorrelation}

\begin{subappendices}
\section{Code for Spatial Stats in R}
\lstinputlisting{"Programs/UP504_SpatStat1.R"}

\renewcommand\bibname{Further reading}

\bibliographystyle{econometrica}
\bibliography{Literature/OverLord}
\IfFileExists{Literature/SpatialStats.bbl}{\chapter{Applied Spatial Data Analysis}\label{chap_SpatialStats}
Quantitative analysis of spatial data has advanced greatly in recent years, largely as the result of the rapidly decreasing cost of computational power, the increased availability of high-resolution data and methodological innovation in spatial econometrics. Indeed, recent analytical and technological advances are increasingly blurring the traditional boundaries in spatial data analysis where various applied methods in Geographic Information Systems (GIS)\index{GIS|see{geographic information systems}}\index{geographic information systems}, geology-oriented remote sensing\index{remote sensing} and spatial statistics\index{spatial statistics} were essentially treated as separate fields. The legacy of this separate treatment of spatial data is reflected by the fact that even the leading textbooks in this area do not focus on more than two of these three sub-fields \citep[e.g.][]{Davis2002a,O'Sullivan2003,Campbell2006}.

Nowhere in the social sciences, perhaps, is this silent revolution felt more than in urban and regional planning. While planners usually have emphasised the application to GIS for land-use and environmental planning purposes, sometimes in addition with a little attention being paid to either remote sensing and spatial statistics, the discipline has never embraced all three sub-fields of quantitative data analysis in a systematic way. 

However, with the rapid deployment of ``Big Data''\index{Big Data} and the emergence of ``policy informatics''\index{policy informatics} as a new paradigm among computational social scientists, planners should increasingly see themselves at the forefront of the application of such integrated spatial data methods.\footnote{See a recent study by McKinsey \& Co. on this \href{http://www.mckinsey.com/Insights/MGI/Research/Technology\_and\_Innovation/Big\_data\_The\_next\_frontier\_for\_innovation}{topic}. See \citet{Barrett2011} for a definition of the term and on recent developments in the area of policy informatics.} \citet{Anselin2012} provides a good overview that assesses the evolution of the way in which spatial data analytical methods have been incorporated into software tools over the past two decades.


\section{Motivations for ``going spatial''}

Spatial autocorrelation\index{spatial autocorrelation|see{autocorrelation}}\index{autocorrelation!spatial} can be loosely defined as ``the coincidence of value similarity with locational similarity'' \citep{Anselin2002}. The problem of spatial dependence then is best viewed as a normal extension of the first law in geography, whereby

\begin{quote}
``everything is related to everything else, but near things are more related than distant things'' \citep[p.236]{Tobler1970}.
\end{quote}

Broadly speaking, there are three main reasons for deploying spatial methods:

\begin{enumerate}[(i)]
\item \emph{The independence assumption is not valid}: The attributes of observation $i$ may influence the attributes of observation
$j$.
\item \emph{Spatial heterogeneity}: The magnitude and direction the the treatment effect may vary across space.
\item \emph{Omitted variable bias}: Some common unobserved or latent variables are shared among spatial neighbours.
\end{enumerate}


Recall that the matrix notation of the non-spatial DGP was defined as

\begin{eqnarray}
\boldsymbol{y} = \boldsymbol{X\beta} + \boldsymbol{e} \qquad \text{with} \qquad \boldsymbol{e}\sim\mathcal{N}(0,\sigma^{2}\boldsymbol{I}).\label{eqn_SS_DGP}
\end{eqnarray}

Similarly, also recall that the classical normal regression model (CNR) of equation~\ref{eqn_SS_DGP} is based on the following assumptions:

\begin{itemize}
\item Observed values at location $i$ are independent of values at location $j$.
\item The regression residuals associated observation $i$ are independent of those associated with observation $j$, i.e. $\mathbb{E}[e_{i}e_{j}]=\mathbb{E}[e_{i}]\mathbb{E}[e_{j}]=0$.
\end{itemize}

\begin{itemize}
\item Include the following literature: \citet{Bivand2008}, \citet{Bivand2008a}, \citet{Corrado2012}, \citet{Kelejian1998}, \citet{McMillen2010a}, \citet{Pinske2010}.
\end{itemize}




\section{Spatial point pattern analysis}
\section{Interpolation and geostatistics}
\section{Areal data and spatial autocorrelation}

\begin{subappendices}
\section{Code for Spatial Stats in R}
\lstinputlisting{"Programs/UP504_SpatStat1.R"}

\renewcommand\bibname{Further reading}

\bibliographystyle{econometrica}
\bibliography{Literature/OverLord}
\IfFileExists{Literature/SpatialStats.bbl}{\input{Literature/SpatialStats.bbl}}{\typeout{Need to run bibtex}}
\end{subappendices} }{\typeout{Need to run bibtex}}
\end{subappendices} }{\typeout{Need to run bibtex}}
\end{subappendices} }{\typeout{Need bibtex on SpatialStats}}
\end{subappendices}
